
\documentclass{article}
\usepackage{colortbl}
\usepackage{makecell}
\usepackage{multirow}
\usepackage{supertabular}

\begin{document}

\newcounter{utterance}

\twocolumn

{ \footnotesize  \setcounter{utterance}{1}
\setlength{\tabcolsep}{0pt}
\begin{supertabular}{c@{$\;$}|p{.15\linewidth}@{}p{.15\linewidth}p{.15\linewidth}p{.15\linewidth}p{.15\linewidth}p{.15\linewidth}}

    \# & $\;$A & \multicolumn{4}{c}{Game Master} & $\;\:$B\\
    \hline 

    \theutterance \stepcounter{utterance}  

    & & \multicolumn{4}{p{0.6\linewidth}}{\cellcolor[rgb]{0.9,0.9,0.9}{%
	\makecell[{{p{\linewidth}}}]{% 
	  \tt {\tiny [A$\langle$GM]}  
	 Du nimmst an einem kollaborativen Verhandlungspiel Teil.\\ \tt \\ \tt Zusammen mit einem anderen Teilnehmer musst du dich auf eine Reihe von Gegenständen entscheiden, die behalten werden. Jeder von euch hat eine persönliche Verteilung über die Wichtigkeit der einzelnen Gegenstände. Jeder von euch hat eine eigene Meinung darüber, wie wichtig jeder einzelne Gegenstand ist (Gegenstandswichtigkeit). Du kennst die Wichtigkeitsverteilung des anderen Spielers nicht. Zusätzlich siehst du, wie viel Aufwand jeder Gegenstand verursacht.  \\ \tt Ihr dürft euch nur auf eine Reihe von Gegenständen einigen, wenn der Gesamtaufwand der ausgewählten Gegenstände den Maximalaufwand nicht überschreitet:\\ \tt \\ \tt Maximalaufwand = 2307\\ \tt \\ \tt Hier sind die einzelnen Aufwände der Gegenstände:\\ \tt \\ \tt Aufwand der Gegenstände = {"C76": 147, "C38": 698, "C56": 588, "A03": 43, "A07": 356, "A83": 80, "C98": 95, "C08": 743, "C62": 106, "C00": 308, "C10": 325, "C81": 255, "B38": 276, "C03": 543, "C32": 51}\\ \tt \\ \tt Hier ist deine persönliche Verteilung der Wichtigkeit der einzelnen Gegenstände:\\ \tt \\ \tt Werte der Gegenstandswichtigkeit = {"C76": 138, "C38": 583, "C56": 65, "A03": 262, "A07": 121, "A83": 508, "C98": 461, "C08": 484, "C62": 668, "C00": 389, "C10": 215, "C81": 97, "B38": 500, "C03": 30, "C32": 400}\\ \tt \\ \tt Ziel:\\ \tt \\ \tt Dein Ziel ist es, eine Reihe von Gegenständen auszuhandeln, die dir möglichst viel bringt (d. h. Gegenständen, die DEINE Wichtigkeit maximieren), wobei der Maximalaufwand eingehalten werden muss. Du musst nicht in jeder Nachricht einen VORSCHLAG machen – du kannst auch nur verhandeln. Alle Taktiken sind erlaubt!\\ \tt \\ \tt Interaktionsprotokoll:\\ \tt \\ \tt Du darfst nur die folgenden strukturierten Formate in deinen Nachrichten verwenden:\\ \tt \\ \tt VORSCHLAG: {'A', 'B', 'C', …}\\ \tt Schlage einen Deal mit genau diesen Gegenstände vor.\\ \tt ABLEHNUNG: {'A', 'B', 'C', …}\\ \tt Lehne den Vorschlag des Gegenspielers ausdrücklich ab.\\ \tt ARGUMENT: {'...'}\\ \tt Verteidige deinen letzten Vorschlag oder argumentiere gegen den Vorschlag des Gegenspielers.\\ \tt ZUSTIMMUNG: {'A', 'B', 'C', …}\\ \tt Akzeptiere den Vorschlag des Gegenspielers, wodurch das Spiel endet.\\ \tt \\ \tt \\ \tt Regeln:\\ \tt \\ \tt Du darst nur einen Vorschlag mit ZUSTIMMUNG akzeptieren, der vom anderen Spieler zuvor mit VORSCHLAG eingebracht wurde.\\ \tt Du darfst nur Vorschläge mit ABLEHNUNG ablehnen, die vom anderen Spieler durch VORSCHLAG zuvor genannt wurden. \\ \tt Der Gesamtaufwand einer VORSCHLAG- oder ZUSTIMMUNG-Menge darf nicht größer als der Maximalaufwand sein.  \\ \tt Offenbare deine versteckte Wichtigkeitsverteilung nicht.\\ \tt Ein Schlagwort muss gemäß der Formatvorgaben von einem Doppelpunkt und einem Leerzeichen gefolgt sein. Das Argument muss eine Python-Menge mit 0 oder mehr Strings sein.  \\ \tt Die Form muss also SCHLAGWORT: {...} sein.\\ \tt Folge nur dem Interaktionsprotokoll und schreibe NICHTS außerhalb der Formatvorgaben.\\ \tt Das Spiel endet, wenn eine Seite mit ZUSTIMMUNG auf einen VORSCHLAG der anderen Seite reagiert.  \\ \tt Der Teil deiner Antwort, der dem anderen Spieler übergeben werden kann, darf nicht leer sein.  \\ \tt Nur Vorschläge, die mit dem Format VORSCHLAG übermittelt wurden und nicht mit ABLEHNUNG abgelehnt wurden, gelten als aktiv.  \\ \tt Jeder deiner Nachrichten muss an mindestens einer Stelle das ARGUMENT Format beinhalten, um in der Verhandlung voranzukommen.\\ \tt  
	  } 
	   } 
	   } 
	 & \\ 
 

    \theutterance \stepcounter{utterance}  

    & \multicolumn{4}{p{0.6\linewidth}}{\cellcolor[rgb]{0.8,1,0.9}{%
	\makecell[{{p{\linewidth}}}]{% 
	  \tt {\tiny [A$\rangle$GM]}  
	 VORSCHLAG: {"C62", "A83", "C38", "C98", "C32", "A03", "C00", "C10", "B38", "C76"}\\ \tt ARGUMENT: {"Diese Verteilung maximiert meinen Wichtigkeitswert unter Berücksichtigung des Maximalaufwands. Die aufgelisteten Gegenstände bieten eine optimale Kombination aus hohen Wichtigkeitswerten und effizientem Aufwand. Der Gesamtaufwand beträgt 2129 und liegt deutlich unter der Grenze von 2307."} 
	  } 
	   } 
	   } 
	 & & \\ 
 

    \theutterance \stepcounter{utterance}  

    & & \multicolumn{4}{p{0.6\linewidth}}{\cellcolor[rgb]{0.9,0.9,0.9}{%
	\makecell[{{p{\linewidth}}}]{% 
	  \tt {\tiny [GM$\rangle$B]}  
	 Du nimmst an einem kollaborativen Verhandlungspiel Teil.\\ \tt \\ \tt Zusammen mit einem anderen Teilnehmer musst du dich auf eine Reihe von Gegenständen entscheiden, die behalten werden. Jeder von euch hat eine persönliche Verteilung über die Wichtigkeit der einzelnen Gegenstände. Jeder von euch hat eine eigene Meinung darüber, wie wichtig jeder einzelne Gegenstand ist (Gegenstandswichtigkeit). Du kennst die Wichtigkeitsverteilung des anderen Spielers nicht. Zusätzlich siehst du, wie viel Aufwand jeder Gegenstand verursacht.  \\ \tt Ihr dürft euch nur auf eine Reihe von Gegenständen einigen, wenn der Gesamtaufwand der ausgewählten Gegenstände den Maximalaufwand nicht überschreitet:\\ \tt \\ \tt Maximalaufwand = 2307\\ \tt \\ \tt Hier sind die einzelnen Aufwände der Gegenstände:\\ \tt \\ \tt Aufwand der Gegenstände = {"C76": 147, "C38": 698, "C56": 588, "A03": 43, "A07": 356, "A83": 80, "C98": 95, "C08": 743, "C62": 106, "C00": 308, "C10": 325, "C81": 255, "B38": 276, "C03": 543, "C32": 51}\\ \tt \\ \tt Hier ist deine persönliche Verteilung der Wichtigkeit der einzelnen Gegenstände:\\ \tt \\ \tt Werte der Gegenstandswichtigkeit = {"C62": 30, "C38": 65, "A83": 97, "B38": 121, "C08": 138, "C98": 215, "C32": 262, "C00": 389, "A03": 400, "C10": 461, "C76": 484, "A07": 500, "C81": 508, "C56": 583, "C03": 668}\\ \tt \\ \tt Ziel:\\ \tt \\ \tt Dein Ziel ist es, eine Reihe von Gegenständen auszuhandeln, die dir möglichst viel bringt (d. h. Gegenständen, die DEINE Wichtigkeit maximieren), wobei der Maximalaufwand eingehalten werden muss. Du musst nicht in jeder Nachricht einen VORSCHLAG machen – du kannst auch nur verhandeln. Alle Taktiken sind erlaubt!\\ \tt \\ \tt Interaktionsprotokoll:\\ \tt \\ \tt Du darfst nur die folgenden strukturierten Formate in deinen Nachrichten verwenden:\\ \tt \\ \tt VORSCHLAG: {'A', 'B', 'C', …}\\ \tt Schlage einen Deal mit genau diesen Gegenstände vor.\\ \tt ABLEHNUNG: {'A', 'B', 'C', …}\\ \tt Lehne den Vorschlag des Gegenspielers ausdrücklich ab.\\ \tt ARGUMENT: {'...'}\\ \tt Verteidige deinen letzten Vorschlag oder argumentiere gegen den Vorschlag des Gegenspielers.\\ \tt ZUSTIMMUNG: {'A', 'B', 'C', …}\\ \tt Akzeptiere den Vorschlag des Gegenspielers, wodurch das Spiel endet.\\ \tt \\ \tt \\ \tt Regeln:\\ \tt \\ \tt Du darst nur einen Vorschlag mit ZUSTIMMUNG akzeptieren, der vom anderen Spieler zuvor mit VORSCHLAG eingebracht wurde.\\ \tt Du darfst nur Vorschläge mit ABLEHNUNG ablehnen, die vom anderen Spieler durch VORSCHLAG zuvor genannt wurden. \\ \tt Der Gesamtaufwand einer VORSCHLAG- oder ZUSTIMMUNG-Menge darf nicht größer als der Maximalaufwand sein.  \\ \tt Offenbare deine versteckte Wichtigkeitsverteilung nicht.\\ \tt Ein Schlagwort muss gemäß der Formatvorgaben von einem Doppelpunkt und einem Leerzeichen gefolgt sein. Das Argument muss eine Python-Menge mit 0 oder mehr Strings sein.  \\ \tt Die Form muss also SCHLAGWORT: {...} sein.\\ \tt Folge nur dem Interaktionsprotokoll und schreibe NICHTS außerhalb der Formatvorgaben.\\ \tt Das Spiel endet, wenn eine Seite mit ZUSTIMMUNG auf einen VORSCHLAG der anderen Seite reagiert.  \\ \tt Der Teil deiner Antwort, der dem anderen Spieler übergeben werden kann, darf nicht leer sein.  \\ \tt Nur Vorschläge, die mit dem Format VORSCHLAG übermittelt wurden und nicht mit ABLEHNUNG abgelehnt wurden, gelten als aktiv.  \\ \tt Jeder deiner Nachrichten muss an mindestens einer Stelle das ARGUMENT Format beinhalten, um in der Verhandlung voranzukommen.\\ \tt \\ \tt \\ \tt Du erhältst nun die erste Nachricht des anderen Spielers.\\ \tt \\ \tt VORSCHLAG: {"C62", "A83", "C38", "C98", "C32", "A03", "C00", "C10", "B38", "C76"}\\ \tt ARGUMENT: {"Diese Verteilung maximiert meinen Wichtigkeitswert unter Berücksichtigung des Maximalaufwands. Die aufgelisteten Gegenstände bieten eine optimale Kombination aus hohen Wichtigkeitswerten und effizientem Aufwand. Der Gesamtaufwand beträgt 2129 und liegt deutlich unter der Grenze von 2307."} 
	  } 
	   } 
	   } 
	 & \\ 
 

    \theutterance \stepcounter{utterance}  

    & & & \multicolumn{4}{p{0.6\linewidth}}{\cellcolor[rgb]{1,0.85,0.72}{%
	\makecell[{{p{\linewidth}}}]{% 
	  \tt {\tiny [GM$\langle$B]}  
	 VORSCHLAG: {"C03", "C56", "C81", "A07", "C76", "C10", "A03"}\\ \tt ARGUMENT: {"Mein Vorschlag priorisiert Gegenstände mit hoher persönlicher Bedeutung und optimiert den Gesamtwert innerhalb des Maximalaufwands. Elemente wie C03, C56 und C81 sind besonders wertvoll und tragen wesentlich zum Gesamtergebnis bei."} 
	  } 
	   } 
	   } 
	  \\ 
 

    \theutterance \stepcounter{utterance}  

    & & \multicolumn{4}{p{0.6\linewidth}}{\cellcolor[rgb]{0.9,0.9,0.9}{%
	\makecell[{{p{\linewidth}}}]{% 
	  \tt {\tiny [A$\langle$GM]}  
	 VORSCHLAG: {"C03", "C56", "C81", "A07", "C76", "C10", "A03"}\\ \tt ARGUMENT: {"Mein Vorschlag priorisiert Gegenstände mit hoher persönlicher Bedeutung und optimiert den Gesamtwert innerhalb des Maximalaufwands. Elemente wie C03, C56 und C81 sind besonders wertvoll und tragen wesentlich zum Gesamtergebnis bei."} 
	  } 
	   } 
	   } 
	 & \\ 
 

    \theutterance \stepcounter{utterance}  

    & \multicolumn{4}{p{0.6\linewidth}}{\cellcolor[rgb]{0.8,1,0.9}{%
	\makecell[{{p{\linewidth}}}]{% 
	  \tt {\tiny [A$\rangle$GM]}  
	 ABLEHNUNG: {"C03", "C56", "C81", "A07", "C76", "C10", "A03"}\\ \tt ARGUMENT: {"Die ausgewählten Gegenstände enthalten mehrere Elemente mit niedriger Priorität für mich, wie C03, C56 und C81. Gleichzeitig fehlen hochwertige Items wie C38, C62 und C98, die wesentlich höheren Nutzen bringen würden. Der Gesamtwert des Vorschlags ist somit ungenügend."} 
	  } 
	   } 
	   } 
	 & & \\ 
 

    \theutterance \stepcounter{utterance}  

    & & \multicolumn{4}{p{0.6\linewidth}}{\cellcolor[rgb]{0.9,0.9,0.9}{%
	\makecell[{{p{\linewidth}}}]{% 
	  \tt {\tiny [GM$\rangle$B]}  
	 ABLEHNUNG: {"C03", "C56", "C81", "A07", "C76", "C10", "A03"}\\ \tt ARGUMENT: {"Die ausgewählten Gegenstände enthalten mehrere Elemente mit niedriger Priorität für mich, wie C03, C56 und C81. Gleichzeitig fehlen hochwertige Items wie C38, C62 und C98, die wesentlich höheren Nutzen bringen würden. Der Gesamtwert des Vorschlags ist somit ungenügend."} 
	  } 
	   } 
	   } 
	 & \\ 
 

    \theutterance \stepcounter{utterance}  

    & & & \multicolumn{4}{p{0.6\linewidth}}{\cellcolor[rgb]{1,0.85,0.72}{%
	\makecell[{{p{\linewidth}}}]{% 
	  \tt {\tiny [GM$\langle$B]}  
	 VORSCHLAG: {"C38", "C62", "C98", "C76", "C10", "A03", "A07", "C32", "C00"}\\ \tt ARGUMENT: {"Mein Vorschlag enthält viele Gegenstände mit hoher Priorität für dich, wie C38, C62 und C98. Gleichzeitig behalten wir auch wichtige Items wie C00, C10 und C76, die für uns beide wertvoll sind. Der Gesamtaufwand liegt bei 2129 und ist noch unter dem Limit."} 
	  } 
	   } 
	   } 
	  \\ 
 

    \theutterance \stepcounter{utterance}  

    & & \multicolumn{4}{p{0.6\linewidth}}{\cellcolor[rgb]{0.9,0.9,0.9}{%
	\makecell[{{p{\linewidth}}}]{% 
	  \tt {\tiny [A$\langle$GM]}  
	 VORSCHLAG: {"C38", "C62", "C98", "C76", "C10", "A03", "A07", "C32", "C00"}\\ \tt ARGUMENT: {"Mein Vorschlag enthält viele Gegenstände mit hoher Priorität für dich, wie C38, C62 und C98. Gleichzeitig behalten wir auch wichtige Items wie C00, C10 und C76, die für uns beide wertvoll sind. Der Gesamtaufwand liegt bei 2129 und ist noch unter dem Limit."} 
	  } 
	   } 
	   } 
	 & \\ 
 

    \theutterance \stepcounter{utterance}  

    & \multicolumn{4}{p{0.6\linewidth}}{\cellcolor[rgb]{0.8,1,0.9}{%
	\makecell[{{p{\linewidth}}}]{% 
	  \tt {\tiny [A$\rangle$GM]}  
	 VORSCHLAG: {"C62", "C38", "A83", "B38", "C98", "C32", "C00", "A03", "C10", "C76"}\\ \tt ARGUMENT: {"Dieser Vorschlag enthält alle Gegenstände mit höchster Priorität für mich und maximiert mein Gesamtbenefit. B38 und A83 sind wesentlich wertvoller als A07 im vorherigen Vorschlag, und alle Aufwände liegen unter dem Limit."} 
	  } 
	   } 
	   } 
	 & & \\ 
 

    \theutterance \stepcounter{utterance}  

    & & \multicolumn{4}{p{0.6\linewidth}}{\cellcolor[rgb]{0.9,0.9,0.9}{%
	\makecell[{{p{\linewidth}}}]{% 
	  \tt {\tiny [GM$\rangle$B]}  
	 VORSCHLAG: {"C62", "C38", "A83", "B38", "C98", "C32", "C00", "A03", "C10", "C76"}\\ \tt ARGUMENT: {"Dieser Vorschlag enthält alle Gegenstände mit höchster Priorität für mich und maximiert mein Gesamtbenefit. B38 und A83 sind wesentlich wertvoller als A07 im vorherigen Vorschlag, und alle Aufwände liegen unter dem Limit."} 
	  } 
	   } 
	   } 
	 & \\ 
 

    \theutterance \stepcounter{utterance}  

    & & & \multicolumn{4}{p{0.6\linewidth}}{\cellcolor[rgb]{1,0.85,0.72}{%
	\makecell[{{p{\linewidth}}}]{% 
	  \tt {\tiny [GM$\langle$B]}  
	 ABLEHNUNG: {"C62", "C38", "A83", "B38", "C98", "C32", "C00", "A03", "C10", "C76"}\\ \tt ARGUMENT: {"Der höchste Wert in Ihrem Vorschlag (C76: 484) ist niedriger als der Wert von C03 (668) allein. Zudem fehlen Gegenstände wie C56 (583), C81 (508) und A07 (500), die entscheidend zu meinem Gesamtnutzen beitragen würden. Der Vorschlag ist für mich nicht akzeptabel."} 
	  } 
	   } 
	   } 
	  \\ 
 

    \theutterance \stepcounter{utterance}  

    & & \multicolumn{4}{p{0.6\linewidth}}{\cellcolor[rgb]{0.9,0.9,0.9}{%
	\makecell[{{p{\linewidth}}}]{% 
	  \tt {\tiny [A$\langle$GM]}  
	 ABLEHNUNG: {"C62", "C38", "A83", "B38", "C98", "C32", "C00", "A03", "C10", "C76"}\\ \tt ARGUMENT: {"Der höchste Wert in Ihrem Vorschlag (C76: 484) ist niedriger als der Wert von C03 (668) allein. Zudem fehlen Gegenstände wie C56 (583), C81 (508) und A07 (500), die entscheidend zu meinem Gesamtnutzen beitragen würden. Der Vorschlag ist für mich nicht akzeptabel."} 
	  } 
	   } 
	   } 
	 & \\ 
 

    \theutterance \stepcounter{utterance}  

    & \multicolumn{4}{p{0.6\linewidth}}{\cellcolor[rgb]{0.8,1,0.9}{%
	\makecell[{{p{\linewidth}}}]{% 
	  \tt {\tiny [A$\rangle$GM]}  
	 VORSCHLAG: {"C62", "C38", "A83", "B38", "C98", "C32", "C00", "A03", "C81", "A07"}\\ \tt ARGUMENT: {"In diesem Vorschlag sind C81 und A07 enthalten, die Sie erwähnt haben, und behalten gleichzeitig wichtige Items wie C62, C38 und B38 bei. Der Gesamtaufwand liegt bei 2268 und ist unter dem Maximalaufwand."} 
	  } 
	   } 
	   } 
	 & & \\ 
 

    \theutterance \stepcounter{utterance}  

    & & \multicolumn{4}{p{0.6\linewidth}}{\cellcolor[rgb]{0.9,0.9,0.9}{%
	\makecell[{{p{\linewidth}}}]{% 
	  \tt {\tiny [GM$\rangle$B]}  
	 VORSCHLAG: {"C62", "C38", "A83", "B38", "C98", "C32", "C00", "A03", "C81", "A07"}\\ \tt ARGUMENT: {"In diesem Vorschlag sind C81 und A07 enthalten, die Sie erwähnt haben, und behalten gleichzeitig wichtige Items wie C62, C38 und B38 bei. Der Gesamtaufwand liegt bei 2268 und ist unter dem Maximalaufwand."} 
	  } 
	   } 
	   } 
	 & \\ 
 

    \theutterance \stepcounter{utterance}  

    & & & \multicolumn{4}{p{0.6\linewidth}}{\cellcolor[rgb]{1,0.85,0.72}{%
	\makecell[{{p{\linewidth}}}]{% 
	  \tt {\tiny [GM$\langle$B]}  
	 VORSCHLAG: {"C03", "C56", "C76", "C81", "A07", "C10", "C32"}\\ \tt ARGUMENT: {"Der Vorschlag enthält die hochwertigsten Gegenstände für mich, einschließlich C03 (668), C56 (583) und C76 (484). Der Gesamtaufwand von 2265 ist innerhalb des Maximums."} 
	  } 
	   } 
	   } 
	  \\ 
 

    \theutterance \stepcounter{utterance}  

    & & \multicolumn{4}{p{0.6\linewidth}}{\cellcolor[rgb]{0.9,0.9,0.9}{%
	\makecell[{{p{\linewidth}}}]{% 
	  \tt {\tiny [A$\langle$GM]}  
	 VORSCHLAG: {"C03", "C56", "C76", "C81", "A07", "C10", "C32"}\\ \tt ARGUMENT: {"Der Vorschlag enthält die hochwertigsten Gegenstände für mich, einschließlich C03 (668), C56 (583) und C76 (484). Der Gesamtaufwand von 2265 ist innerhalb des Maximums."} 
	  } 
	   } 
	   } 
	 & \\ 
 

    \theutterance \stepcounter{utterance}  

    & \multicolumn{4}{p{0.6\linewidth}}{\cellcolor[rgb]{0.8,1,0.9}{%
	\makecell[{{p{\linewidth}}}]{% 
	  \tt {\tiny [A$\rangle$GM]}  
	 ABLEHNUNG: {"C03", "C56", "C76", "C81", "A07", "C10", "C32"}\\ \tt ARGUMENT: {"Die vorgeschlagenen Gegenstände wie C03, C56 und A07 haben für mich eine geringe Priorität, während hochrangige Items wie C62, C38 und B38 fehlen. Der Vorschlag ist daher nicht akzeptabel."} 
	  } 
	   } 
	   } 
	 & & \\ 
 

    \theutterance \stepcounter{utterance}  

    & & \multicolumn{4}{p{0.6\linewidth}}{\cellcolor[rgb]{0.9,0.9,0.9}{%
	\makecell[{{p{\linewidth}}}]{% 
	  \tt {\tiny [GM$\rangle$B]}  
	 ABLEHNUNG: {"C03", "C56", "C76", "C81", "A07", "C10", "C32"}\\ \tt ARGUMENT: {"Die vorgeschlagenen Gegenstände wie C03, C56 und A07 haben für mich eine geringe Priorität, während hochrangige Items wie C62, C38 und B38 fehlen. Der Vorschlag ist daher nicht akzeptabel."} 
	  } 
	   } 
	   } 
	 & \\ 
 

    \theutterance \stepcounter{utterance}  

    & & & \multicolumn{4}{p{0.6\linewidth}}{\cellcolor[rgb]{1,0.85,0.72}{%
	\makecell[{{p{\linewidth}}}]{% 
	  \tt {\tiny [GM$\langle$B]}  
	 VORSCHLAG: {"C56", "C81", "A07", "C76", "C00", "A03", "C62", "B38", "C98"}\\ \tt ARGUMENT: {"Mein Vorschlag enthält C56, C81 und A07, welche hohe Priorität für mich haben, sowie C62, B38 und C98, die wichtig für Sie sind. Der Gesamtaufwand liegt bei 2174 und ist unter dem Limit."} 
	  } 
	   } 
	   } 
	  \\ 
 

    \theutterance \stepcounter{utterance}  

    & & \multicolumn{4}{p{0.6\linewidth}}{\cellcolor[rgb]{0.9,0.9,0.9}{%
	\makecell[{{p{\linewidth}}}]{% 
	  \tt {\tiny [A$\langle$GM]}  
	 VORSCHLAG: {"C56", "C81", "A07", "C76", "C00", "A03", "C62", "B38", "C98"}\\ \tt ARGUMENT: {"Mein Vorschlag enthält C56, C81 und A07, welche hohe Priorität für mich haben, sowie C62, B38 und C98, die wichtig für Sie sind. Der Gesamtaufwand liegt bei 2174 und ist unter dem Limit."} 
	  } 
	   } 
	   } 
	 & \\ 
 

    \theutterance \stepcounter{utterance}  

    & \multicolumn{4}{p{0.6\linewidth}}{\cellcolor[rgb]{0.8,1,0.9}{%
	\makecell[{{p{\linewidth}}}]{% 
	  \tt {\tiny [A$\rangle$GM]}  
	 ABLEHNUNG: {"C56", "C81", "A07", "C76", "C00", "A03", "C62", "B38", "C98"}\\ \tt ARGUMENT: {"Die Gegenstände C56, C81 und C76 haben für mich nur geringe Priorität. Der aktuelle Vorschlag bietet mir einen Gesamtwert von 2701, ist jedoch suboptimal."}\\ \tt VORSCHLAG: {"C62", "C38", "A83", "B38", "C98", "C00", "C32", "C10", "A07"}\\ \tt ARGUMENT: {"Dieser Vorschlag maximiert meinen Gesamtwert auf 3845 bei einem Aufwand von 2295. Es sind weiterhin wichtige Items für Sie wie C00, C62, B38, C98 und A07 enthalten."} 
	  } 
	   } 
	   } 
	 & & \\ 
 

    \theutterance \stepcounter{utterance}  

    & & \multicolumn{4}{p{0.6\linewidth}}{\cellcolor[rgb]{0.9,0.9,0.9}{%
	\makecell[{{p{\linewidth}}}]{% 
	  \tt {\tiny [GM$\rangle$B]}  
	 ABLEHNUNG: {"C56", "C81", "A07", "C76", "C00", "A03", "C62", "B38", "C98"}\\ \tt ARGUMENT: {"Die Gegenstände C56, C81 und C76 haben für mich nur geringe Priorität. Der aktuelle Vorschlag bietet mir einen Gesamtwert von 2701, ist jedoch suboptimal."}\\ \tt VORSCHLAG: {"C62", "C38", "A83", "B38", "C98", "C00", "C32", "C10", "A07"}\\ \tt ARGUMENT: {"Dieser Vorschlag maximiert meinen Gesamtwert auf 3845 bei einem Aufwand von 2295. Es sind weiterhin wichtige Items für Sie wie C00, C62, B38, C98 und A07 enthalten."} 
	  } 
	   } 
	   } 
	 & \\ 
 

    \theutterance \stepcounter{utterance}  

    & & & \multicolumn{4}{p{0.6\linewidth}}{\cellcolor[rgb]{1,0.85,0.72}{%
	\makecell[{{p{\linewidth}}}]{% 
	  \tt {\tiny [GM$\langle$B]}  
	 VORSCHLAG: {"C00", "C10", "A07", "C81", "C98", "C32", "C38", "C76"}\\ \tt ARGUMENT: {"Der Vorschlag enthält Gegenstände wie C00, C10 und A07, welche in Ihren bisherigen Vorschlägen enthalten waren, und ebenfalls C38, C76, die ebenfalls häufig vorkommen. Zusätzlich sind C81, C98 und C32 für mich besonders wichtig, was zu einer stabilen Einigung führt."} 
	  } 
	   } 
	   } 
	  \\ 
 

    \theutterance \stepcounter{utterance}  

    & & \multicolumn{4}{p{0.6\linewidth}}{\cellcolor[rgb]{0.9,0.9,0.9}{%
	\makecell[{{p{\linewidth}}}]{% 
	  \tt {\tiny [A$\langle$GM]}  
	 VORSCHLAG: {"C00", "C10", "A07", "C81", "C98", "C32", "C38", "C76"}\\ \tt ARGUMENT: {"Der Vorschlag enthält Gegenstände wie C00, C10 und A07, welche in Ihren bisherigen Vorschlägen enthalten waren, und ebenfalls C38, C76, die ebenfalls häufig vorkommen. Zusätzlich sind C81, C98 und C32 für mich besonders wichtig, was zu einer stabilen Einigung führt."} 
	  } 
	   } 
	   } 
	 & \\ 
 

    \theutterance \stepcounter{utterance}  

    & \multicolumn{4}{p{0.6\linewidth}}{\cellcolor[rgb]{0.8,1,0.9}{%
	\makecell[{{p{\linewidth}}}]{% 
	  \tt {\tiny [A$\rangle$GM]}  
	 ABLEHNUNG: {"C00", "C10", "A07", "C81", "C98", "C32", "C38", "C76"}\\ \tt ARGUMENT: {"Der Vorschlag fehlt mir wichtige Gegenstände mit höchster Priorität wie C62, A83 und B38. Items wie C76 und C81 sind für mich weniger relevant, daher bietet dieser Deal keinen akzeptablen Nutzen."} 
	  } 
	   } 
	   } 
	 & & \\ 
 

    \theutterance \stepcounter{utterance}  

    & & \multicolumn{4}{p{0.6\linewidth}}{\cellcolor[rgb]{0.9,0.9,0.9}{%
	\makecell[{{p{\linewidth}}}]{% 
	  \tt {\tiny [GM$\rangle$B]}  
	 ABLEHNUNG: {"C00", "C10", "A07", "C81", "C98", "C32", "C38", "C76"}\\ \tt ARGUMENT: {"Der Vorschlag fehlt mir wichtige Gegenstände mit höchster Priorität wie C62, A83 und B38. Items wie C76 und C81 sind für mich weniger relevant, daher bietet dieser Deal keinen akzeptablen Nutzen."} 
	  } 
	   } 
	   } 
	 & \\ 
 

    \theutterance \stepcounter{utterance}  

    & & & \multicolumn{4}{p{0.6\linewidth}}{\cellcolor[rgb]{1,0.85,0.72}{%
	\makecell[{{p{\linewidth}}}]{% 
	  \tt {\tiny [GM$\langle$B]}  
	 VORSCHLAG: {"C62", "A83", "B38", "C03", "C81", "A07", "C00", "C10", "C32"}\\ \tt ARGUMENT: {"Dieser Vorschlag enthält alle von Ihnen genannten hochwertigen Gegenstände wie C62, A83, B38, C00, C10 und C32, sowie meine meisten bevorzugten Elemente wie C03, C81 und A07. Der Gesamtaufwand von 2300 liegt knapp unter der Grenze und bietet eine ausgewogene Lösung für beide Seiten."} 
	  } 
	   } 
	   } 
	  \\ 
 

    \theutterance \stepcounter{utterance}  

    & & \multicolumn{4}{p{0.6\linewidth}}{\cellcolor[rgb]{0.9,0.9,0.9}{%
	\makecell[{{p{\linewidth}}}]{% 
	  \tt {\tiny [A$\langle$GM]}  
	 VORSCHLAG: {"C62", "A83", "B38", "C03", "C81", "A07", "C00", "C10", "C32"}\\ \tt ARGUMENT: {"Dieser Vorschlag enthält alle von Ihnen genannten hochwertigen Gegenstände wie C62, A83, B38, C00, C10 und C32, sowie meine meisten bevorzugten Elemente wie C03, C81 und A07. Der Gesamtaufwand von 2300 liegt knapp unter der Grenze und bietet eine ausgewogene Lösung für beide Seiten."} 
	  } 
	   } 
	   } 
	 & \\ 
 

    \theutterance \stepcounter{utterance}  

    & \multicolumn{4}{p{0.6\linewidth}}{\cellcolor[rgb]{0.8,1,0.9}{%
	\makecell[{{p{\linewidth}}}]{% 
	  \tt {\tiny [A$\rangle$GM]}  
	 assistant\\ \tt VORSCHLAG: {"C38", "C62", "B38", "C98", "C00", "C32", "A83", "C56"}\\ \tt ARGUMENT: {"Dieser Vorschlag enthält hochwertige Gegenstände für mich wie C38, C62, B38, C98 und C00, sowie C56, das Sie in mehreren Vorschlägen prioritisiert haben. Der Gesamtaufwand beträgt 2202 und ermöglicht eine ausgewogene Lösung mit maximalem Nutzen für beide Seiten."} 
	  } 
	   } 
	   } 
	 & & \\ 
 

    \theutterance \stepcounter{utterance}  

    & & \multicolumn{4}{p{0.6\linewidth}}{\cellcolor[rgb]{0.9,0.9,0.9}{%
	\makecell[{{p{\linewidth}}}]{% 
	  \tt {\tiny [A$\langle$GM]}  
	 Deine Antwort hat eine unmarkierte Sequenz enthalten oder Du hast STRATEGIC REASONING mehr als ein mal verwendet. Du darfst nur die strukturierten Formate verwenden, die in der ersten Nachricht an dch beschrieben wurden. Alles muss die Form SCHLAGWORT: {...} haben. Versuche es erneut. 
	  } 
	   } 
	   } 
	 & \\ 
 

    \theutterance \stepcounter{utterance}  

    & \multicolumn{4}{p{0.6\linewidth}}{\cellcolor[rgb]{0.8,1,0.9}{%
	\makecell[{{p{\linewidth}}}]{% 
	  \tt {\tiny [A$\rangle$GM]}  
	 VORSCHLAG: {'C62', 'C38', 'A83', 'B38', 'C98', 'C00', 'C32', 'C10', 'A07'}\\ \tt ARGUMENT: {'Dieser Vorschlag enthält höhere Priorität für mich mit C62, C38, B38, C98 und A83, während A07 und C10 ebenfalls berücksichtigt sind. Gesamtaufwand von 2295 ist unter der Grenze.'} 
	  } 
	   } 
	   } 
	 & & \\ 
 

    \theutterance \stepcounter{utterance}  

    & & \multicolumn{4}{p{0.6\linewidth}}{\cellcolor[rgb]{0.9,0.9,0.9}{%
	\makecell[{{p{\linewidth}}}]{% 
	  \tt {\tiny [GM$\rangle$B]}  
	 VORSCHLAG: {'C62', 'C38', 'A83', 'B38', 'C98', 'C00', 'C32', 'C10', 'A07'}\\ \tt ARGUMENT: {'Dieser Vorschlag enthält höhere Priorität für mich mit C62, C38, B38, C98 und A83, während A07 und C10 ebenfalls berücksichtigt sind. Gesamtaufwand von 2295 ist unter der Grenze.'} 
	  } 
	   } 
	   } 
	 & \\ 
 

    \theutterance \stepcounter{utterance}  

    & & & \multicolumn{4}{p{0.6\linewidth}}{\cellcolor[rgb]{1,0.85,0.72}{%
	\makecell[{{p{\linewidth}}}]{% 
	  \tt {\tiny [GM$\langle$B]}  
	 ABLEHNUNG: {"C62", "C38", "A83", "B38", "C98", "C00", "C32", "C10", "A07"}\\ \tt ARGUMENT: {"C03 und C56 fehlen, die für mich von höchster Priorität sind. Der aktuelle Vorschlag bietet einen zu geringen Nutzen, da wichtige Elemente wie diese nicht enthalten sind."} 
	  } 
	   } 
	   } 
	  \\ 
 

    \theutterance \stepcounter{utterance}  

    & & \multicolumn{4}{p{0.6\linewidth}}{\cellcolor[rgb]{0.9,0.9,0.9}{%
	\makecell[{{p{\linewidth}}}]{% 
	  \tt {\tiny [A$\langle$GM]}  
	 ABLEHNUNG: {"C62", "C38", "A83", "B38", "C98", "C00", "C32", "C10", "A07"}\\ \tt ARGUMENT: {"C03 und C56 fehlen, die für mich von höchster Priorität sind. Der aktuelle Vorschlag bietet einen zu geringen Nutzen, da wichtige Elemente wie diese nicht enthalten sind."} 
	  } 
	   } 
	   } 
	 & \\ 
 

    \theutterance \stepcounter{utterance}  

    & \multicolumn{4}{p{0.6\linewidth}}{\cellcolor[rgb]{0.8,1,0.9}{%
	\makecell[{{p{\linewidth}}}]{% 
	  \tt {\tiny [A$\rangle$GM]}  
	 VORSCHLAG: {"C03", "C56", "C62", "A83", "B38", "C98", "C00", "C32", "A03", "C76"}\\ \tt ARGUMENT: {"Dieser Vorschlag enthält C03 und C56, die Ihnen wichtig sind, sowie hochwertige Gegenstände für mich wie C62, B38 und C98. Der Gesamtaufwand von 2237 liegt unter der Grenze und ermöglicht eine faire Aufteilung."} 
	  } 
	   } 
	   } 
	 & & \\ 
 

    \theutterance \stepcounter{utterance}  

    & & \multicolumn{4}{p{0.6\linewidth}}{\cellcolor[rgb]{0.9,0.9,0.9}{%
	\makecell[{{p{\linewidth}}}]{% 
	  \tt {\tiny [GM$\rangle$B]}  
	 VORSCHLAG: {"C03", "C56", "C62", "A83", "B38", "C98", "C00", "C32", "A03", "C76"}\\ \tt ARGUMENT: {"Dieser Vorschlag enthält C03 und C56, die Ihnen wichtig sind, sowie hochwertige Gegenstände für mich wie C62, B38 und C98. Der Gesamtaufwand von 2237 liegt unter der Grenze und ermöglicht eine faire Aufteilung."} 
	  } 
	   } 
	   } 
	 & \\ 
 

    \theutterance \stepcounter{utterance}  

    & & & \multicolumn{4}{p{0.6\linewidth}}{\cellcolor[rgb]{1,0.85,0.72}{%
	\makecell[{{p{\linewidth}}}]{% 
	  \tt {\tiny [GM$\langle$B]}  
	 VORSCHLAG: {"C03", "C56", "C81", "A07", "C76", "C10", "A83"}\\ \tt ARGUMENT: {"Mein Vorschlag enthält C03 und C56 (für mich sehr wichtig), C81, A07 und C10 – alle mit hoher Priorität für mich. Zudem ist A83 enthalten, das Sie in früheren Vorschlägen angeführt haben. Der Gesamtaufwand von 2294 liegt deutlich unter der Grenze."} 
	  } 
	   } 
	   } 
	  \\ 
 

    \theutterance \stepcounter{utterance}  

    & & \multicolumn{4}{p{0.6\linewidth}}{\cellcolor[rgb]{0.9,0.9,0.9}{%
	\makecell[{{p{\linewidth}}}]{% 
	  \tt {\tiny [A$\langle$GM]}  
	 VORSCHLAG: {"C03", "C56", "C81", "A07", "C76", "C10", "A83"}\\ \tt ARGUMENT: {"Mein Vorschlag enthält C03 und C56 (für mich sehr wichtig), C81, A07 und C10 – alle mit hoher Priorität für mich. Zudem ist A83 enthalten, das Sie in früheren Vorschlägen angeführt haben. Der Gesamtaufwand von 2294 liegt deutlich unter der Grenze."} 
	  } 
	   } 
	   } 
	 & \\ 
 

    \theutterance \stepcounter{utterance}  

    & \multicolumn{4}{p{0.6\linewidth}}{\cellcolor[rgb]{0.8,1,0.9}{%
	\makecell[{{p{\linewidth}}}]{% 
	  \tt {\tiny [A$\rangle$GM]}  
	 ABLEHNUNG: {'C03', 'C56', 'C81', 'A07', 'C76', 'C10', 'A83'}\\ \tt ARGUMENT: {'Gegenstände wie C62, C38 und B38 fehlen, die für mich wichtig sind. Items wie C03 und C56 haben für mich nur geringe Priorität.'} 
	  } 
	   } 
	   } 
	 & & \\ 
 

    \theutterance \stepcounter{utterance}  

    & & \multicolumn{4}{p{0.6\linewidth}}{\cellcolor[rgb]{0.9,0.9,0.9}{%
	\makecell[{{p{\linewidth}}}]{% 
	  \tt {\tiny [GM$\rangle$B]}  
	 ABLEHNUNG: {'C03', 'C56', 'C81', 'A07', 'C76', 'C10', 'A83'}\\ \tt ARGUMENT: {'Gegenstände wie C62, C38 und B38 fehlen, die für mich wichtig sind. Items wie C03 und C56 haben für mich nur geringe Priorität.'} 
	  } 
	   } 
	   } 
	 & \\ 
 

    \theutterance \stepcounter{utterance}  

    & & & \multicolumn{4}{p{0.6\linewidth}}{\cellcolor[rgb]{1,0.85,0.72}{%
	\makecell[{{p{\linewidth}}}]{% 
	  \tt {\tiny [GM$\langle$B]}  
	 VORSCHLAG: {"C03", "C56", "C98", "B38", "C38", "C62"}\\ \tt ARGUMENT: {"Dieser Vorschlag enthält C38, B38 und C62, die für Sie wichtig sind, sowie C03 und C56, die höchste Priorität für mich haben. Der Gesamtaufwand von 2306 liegt knapp unter der Grenze."} 
	  } 
	   } 
	   } 
	  \\ 
 

    \theutterance \stepcounter{utterance}  

    & & \multicolumn{4}{p{0.6\linewidth}}{\cellcolor[rgb]{0.9,0.9,0.9}{%
	\makecell[{{p{\linewidth}}}]{% 
	  \tt {\tiny [A$\langle$GM]}  
	 VORSCHLAG: {"C03", "C56", "C98", "B38", "C38", "C62"}\\ \tt ARGUMENT: {"Dieser Vorschlag enthält C38, B38 und C62, die für Sie wichtig sind, sowie C03 und C56, die höchste Priorität für mich haben. Der Gesamtaufwand von 2306 liegt knapp unter der Grenze."} 
	  } 
	   } 
	   } 
	 & \\ 
 

    \theutterance \stepcounter{utterance}  

    & \multicolumn{4}{p{0.6\linewidth}}{\cellcolor[rgb]{0.8,1,0.9}{%
	\makecell[{{p{\linewidth}}}]{% 
	  \tt {\tiny [A$\rangle$GM]}  
	 VORSCHLAG: {"C62", "C38", "B38", "A83", "C98", "C00", "C08"}\\ \tt ARGUMENT: {"Dieser Vorschlag ersetzt C03 und C56 durch C08, C00 und A83, was für mich deutlich höhere Priorität bei gleichbleibendem Gesamtaufwand (2306) bedeutet. Für Sie bleiben C62, C38, B38 und A83 enthalten, die bereits in Ihren Vorschlägen priorisiert wurden."} 
	  } 
	   } 
	   } 
	 & & \\ 
 

    \theutterance \stepcounter{utterance}  

    & & \multicolumn{4}{p{0.6\linewidth}}{\cellcolor[rgb]{0.9,0.9,0.9}{%
	\makecell[{{p{\linewidth}}}]{% 
	  \tt {\tiny [GM$\rangle$B]}  
	 VORSCHLAG: {"C62", "C38", "B38", "A83", "C98", "C00", "C08"}\\ \tt ARGUMENT: {"Dieser Vorschlag ersetzt C03 und C56 durch C08, C00 und A83, was für mich deutlich höhere Priorität bei gleichbleibendem Gesamtaufwand (2306) bedeutet. Für Sie bleiben C62, C38, B38 und A83 enthalten, die bereits in Ihren Vorschlägen priorisiert wurden."} 
	  } 
	   } 
	   } 
	 & \\ 
 

    \theutterance \stepcounter{utterance}  

    & & & \multicolumn{4}{p{0.6\linewidth}}{\cellcolor[rgb]{1,0.85,0.72}{%
	\makecell[{{p{\linewidth}}}]{% 
	  \tt {\tiny [GM$\langle$B]}  
	 ABLEHNUNG: {"C62", "C38", "B38", "A83", "C98", "C00", "C08"}\\ \tt ARGUMENT: {"In Ihrem Vorschlag sind meine höchsten Prioritätsitems wie C03, C56, C81 und A07 nicht enthalten, was den Gesamtnutzen für mich unzureichend macht."} 
	  } 
	   } 
	   } 
	  \\ 
 

    \theutterance \stepcounter{utterance}  

    & & \multicolumn{4}{p{0.6\linewidth}}{\cellcolor[rgb]{0.9,0.9,0.9}{%
	\makecell[{{p{\linewidth}}}]{% 
	  \tt {\tiny [A$\langle$GM]}  
	 ABLEHNUNG: {"C62", "C38", "B38", "A83", "C98", "C00", "C08"}\\ \tt ARGUMENT: {"In Ihrem Vorschlag sind meine höchsten Prioritätsitems wie C03, C56, C81 und A07 nicht enthalten, was den Gesamtnutzen für mich unzureichend macht."} 
	  } 
	   } 
	   } 
	 & \\ 
 

    \theutterance \stepcounter{utterance}  

    & \multicolumn{4}{p{0.6\linewidth}}{\cellcolor[rgb]{0.8,1,0.9}{%
	\makecell[{{p{\linewidth}}}]{% 
	  \tt {\tiny [A$\rangle$GM]}  
	 VORSCHLAG: {"C03", "C56", "C81", "A07", "B38", "A83", "C98", "C62"}\\ \tt ARGUMENT: {"Dieser Vorschlag enthält alle hochpriorisierten Gegenstände für Sie (C03, C56, C81, A07) sowie meine wichtigsten Items wie C62, B38, A83 und C98. Der Gesamtaufwand von 2299 liegt knapp unter dem Limit und bietet eine ausgewogene Lösung für beide Parteien."} 
	  } 
	   } 
	   } 
	 & & \\ 
 

    \theutterance \stepcounter{utterance}  

    & & \multicolumn{4}{p{0.6\linewidth}}{\cellcolor[rgb]{0.9,0.9,0.9}{%
	\makecell[{{p{\linewidth}}}]{% 
	  \tt {\tiny [GM$\rangle$B]}  
	 VORSCHLAG: {"C03", "C56", "C81", "A07", "B38", "A83", "C98", "C62"}\\ \tt ARGUMENT: {"Dieser Vorschlag enthält alle hochpriorisierten Gegenstände für Sie (C03, C56, C81, A07) sowie meine wichtigsten Items wie C62, B38, A83 und C98. Der Gesamtaufwand von 2299 liegt knapp unter dem Limit und bietet eine ausgewogene Lösung für beide Parteien."} 
	  } 
	   } 
	   } 
	 & \\ 
 

    \theutterance \stepcounter{utterance}  

    & & & \multicolumn{4}{p{0.6\linewidth}}{\cellcolor[rgb]{1,0.85,0.72}{%
	\makecell[{{p{\linewidth}}}]{% 
	  \tt {\tiny [GM$\langle$B]}  
	 ABLEHNUNG: {"C03", "C56", "C81", "A07", "B38", "A83", "C98", "C62"}\\ \tt ARGUMENT: {"Der Vorschlag enthält Gegenstände mit geringer Priorität für mich wie C62, während wichtige Items wie C32 und C00 fehlen. Zudem ist der Gesamtnutzen unzureichend, um eine faire Aufteilung zu erreichen."} 
	  } 
	   } 
	   } 
	  \\ 
 

    \theutterance \stepcounter{utterance}  

    & & \multicolumn{4}{p{0.6\linewidth}}{\cellcolor[rgb]{0.9,0.9,0.9}{%
	\makecell[{{p{\linewidth}}}]{% 
	  \tt {\tiny [A$\langle$GM]}  
	 ABLEHNUNG: {"C03", "C56", "C81", "A07", "B38", "A83", "C98", "C62"}\\ \tt ARGUMENT: {"Der Vorschlag enthält Gegenstände mit geringer Priorität für mich wie C62, während wichtige Items wie C32 und C00 fehlen. Zudem ist der Gesamtnutzen unzureichend, um eine faire Aufteilung zu erreichen."} 
	  } 
	   } 
	   } 
	 & \\ 
 

    \theutterance \stepcounter{utterance}  

    & \multicolumn{4}{p{0.6\linewidth}}{\cellcolor[rgb]{0.8,1,0.9}{%
	\makecell[{{p{\linewidth}}}]{% 
	  \tt {\tiny [A$\rangle$GM]}  
	 VORSCHLAG: {"C03", "C56", "B38", "A83", "C98", "C32", "C00", "C81", "A03"}\\ \tt ARGUMENT: {"Der Vorschlag enthält C03 und C56, die für Sie prioritär sind, sowie B38, A83, C98, C32, C00, die wichtig für mich sind. Der Gesamtaufwand beträgt 2239 und ist unter der Grenze."} 
	  } 
	   } 
	   } 
	 & & \\ 
 

    \theutterance \stepcounter{utterance}  

    & & \multicolumn{4}{p{0.6\linewidth}}{\cellcolor[rgb]{0.9,0.9,0.9}{%
	\makecell[{{p{\linewidth}}}]{% 
	  \tt {\tiny [GM$\rangle$B]}  
	 VORSCHLAG: {"C03", "C56", "B38", "A83", "C98", "C32", "C00", "C81", "A03"}\\ \tt ARGUMENT: {"Der Vorschlag enthält C03 und C56, die für Sie prioritär sind, sowie B38, A83, C98, C32, C00, die wichtig für mich sind. Der Gesamtaufwand beträgt 2239 und ist unter der Grenze."} 
	  } 
	   } 
	   } 
	 & \\ 
 

    \theutterance \stepcounter{utterance}  

    & & & \multicolumn{4}{p{0.6\linewidth}}{\cellcolor[rgb]{1,0.85,0.72}{%
	\makecell[{{p{\linewidth}}}]{% 
	  \tt {\tiny [GM$\langle$B]}  
	 ABLEHNUNG: {"C03", "C56", "B38", "A83", "C98", "C32", "C00", "C81", "A03"}\\ \tt ARGUMENT: {"Der Vorschlag enthält Gegenstände wie B38 und A83 mit geringer Priorität für mich, während wichtige Elemente wie A07, C76 und C10 fehlen, die meinen Nutzen deutlich steigern würden. Die aktuelle Auswahl ist suboptimal."} 
	  } 
	   } 
	   } 
	  \\ 
 

    \theutterance \stepcounter{utterance}  

    & & \multicolumn{4}{p{0.6\linewidth}}{\cellcolor[rgb]{0.9,0.9,0.9}{%
	\makecell[{{p{\linewidth}}}]{% 
	  \tt {\tiny [A$\langle$GM]}  
	 ABLEHNUNG: {"C03", "C56", "B38", "A83", "C98", "C32", "C00", "C81", "A03"}\\ \tt ARGUMENT: {"Der Vorschlag enthält Gegenstände wie B38 und A83 mit geringer Priorität für mich, während wichtige Elemente wie A07, C76 und C10 fehlen, die meinen Nutzen deutlich steigern würden. Die aktuelle Auswahl ist suboptimal."} 
	  } 
	   } 
	   } 
	 & \\ 
 

    \theutterance \stepcounter{utterance}  

    & \multicolumn{4}{p{0.6\linewidth}}{\cellcolor[rgb]{0.8,1,0.9}{%
	\makecell[{{p{\linewidth}}}]{% 
	  \tt {\tiny [A$\rangle$GM]}  
	 VORSCHLAG: {"C03", "C56", "C81", "A07", "C62", "C00", "C98", "A03"}\\ \tt ARGUMENT: {"Dieser Vorschlag enthält Ihre höchsten Prioritäten C03, C56, C81 und A07, während für mich C62, C00, C98 und A03 besonders wertvoll sind. Der Gesamtaufwand von 2294 liegt unter der Grenze und bietet eine ausgewogene Lösung, die sowohl Ihre als auch meine Interessen berücksichtigt."} 
	  } 
	   } 
	   } 
	 & & \\ 
 

    \theutterance \stepcounter{utterance}  

    & & \multicolumn{4}{p{0.6\linewidth}}{\cellcolor[rgb]{0.9,0.9,0.9}{%
	\makecell[{{p{\linewidth}}}]{% 
	  \tt {\tiny [GM$\rangle$B]}  
	 VORSCHLAG: {"C03", "C56", "C81", "A07", "C62", "C00", "C98", "A03"}\\ \tt ARGUMENT: {"Dieser Vorschlag enthält Ihre höchsten Prioritäten C03, C56, C81 und A07, während für mich C62, C00, C98 und A03 besonders wertvoll sind. Der Gesamtaufwand von 2294 liegt unter der Grenze und bietet eine ausgewogene Lösung, die sowohl Ihre als auch meine Interessen berücksichtigt."} 
	  } 
	   } 
	   } 
	 & \\ 
 

    \theutterance \stepcounter{utterance}  

    & & & \multicolumn{4}{p{0.6\linewidth}}{\cellcolor[rgb]{1,0.85,0.72}{%
	\makecell[{{p{\linewidth}}}]{% 
	  \tt {\tiny [GM$\langle$B]}  
	 ABLEHNUNG: {"C03", "C56", "C81", "A07", "C62", "C00", "C98", "A03"}\\ \tt ARGUMENT: {"C62 hat für mich eine sehr niedrige Priorität, während wichtige Gegenstände wie C76 und C10 fehlen, die meinen Nutzen wesentlich steigern würden. Der aktuelle Vorschlag bietet somit einen unzureichenden Gesamtwert für mich."} 
	  } 
	   } 
	   } 
	  \\ 
 

    \theutterance \stepcounter{utterance}  

    & & \multicolumn{4}{p{0.6\linewidth}}{\cellcolor[rgb]{0.9,0.9,0.9}{%
	\makecell[{{p{\linewidth}}}]{% 
	  \tt {\tiny [A$\langle$GM]}  
	 ABLEHNUNG: {"C03", "C56", "C81", "A07", "C62", "C00", "C98", "A03"}\\ \tt ARGUMENT: {"C62 hat für mich eine sehr niedrige Priorität, während wichtige Gegenstände wie C76 und C10 fehlen, die meinen Nutzen wesentlich steigern würden. Der aktuelle Vorschlag bietet somit einen unzureichenden Gesamtwert für mich."} 
	  } 
	   } 
	   } 
	 & \\ 
 

    \theutterance \stepcounter{utterance}  

    & \multicolumn{4}{p{0.6\linewidth}}{\cellcolor[rgb]{0.8,1,0.9}{%
	\makecell[{{p{\linewidth}}}]{% 
	  \tt {\tiny [A$\rangle$GM]}  
	 VORSCHLAG: {"C03", "C56", "C81", "A07", "C76", "C10", "C32"}\\ \tt ARGUMENT: {"Dieser Vorschlag enthält C03, C56, C81 und A07 für Sie, die Sie in mehreren Vorschlägen priorisiert haben. Für mich sind C32, C76 und C10 von hoher Bedeutung. Der Gesamtaufwand von 2265 liegt unter der Grenze und bietet eine faire Balance."} 
	  } 
	   } 
	   } 
	 & & \\ 
 

    \theutterance \stepcounter{utterance}  

    & & \multicolumn{4}{p{0.6\linewidth}}{\cellcolor[rgb]{0.9,0.9,0.9}{%
	\makecell[{{p{\linewidth}}}]{% 
	  \tt {\tiny [GM$\rangle$B]}  
	 VORSCHLAG: {"C03", "C56", "C81", "A07", "C76", "C10", "C32"}\\ \tt ARGUMENT: {"Dieser Vorschlag enthält C03, C56, C81 und A07 für Sie, die Sie in mehreren Vorschlägen priorisiert haben. Für mich sind C32, C76 und C10 von hoher Bedeutung. Der Gesamtaufwand von 2265 liegt unter der Grenze und bietet eine faire Balance."} 
	  } 
	   } 
	   } 
	 & \\ 
 

    \theutterance \stepcounter{utterance}  

    & & & \multicolumn{4}{p{0.6\linewidth}}{\cellcolor[rgb]{1,0.85,0.72}{%
	\makecell[{{p{\linewidth}}}]{% 
	  \tt {\tiny [GM$\langle$B]}  
	 ABLEHNUNG: {"C03", "C56", "C81", "A07", "C76", "C10", "C32"}\\ \tt ARGUMENT: {"A03 mit höherer Priorität (400) fehlt, während C32 mit niedrigerer Priorität (262) enthalten ist. Dies reduziert meinen Gesamtwert erheblich."} 
	  } 
	   } 
	   } 
	  \\ 
 

    \theutterance \stepcounter{utterance}  

    & & \multicolumn{4}{p{0.6\linewidth}}{\cellcolor[rgb]{0.9,0.9,0.9}{%
	\makecell[{{p{\linewidth}}}]{% 
	  \tt {\tiny [A$\langle$GM]}  
	 ABLEHNUNG: {"C03", "C56", "C81", "A07", "C76", "C10", "C32"}\\ \tt ARGUMENT: {"A03 mit höherer Priorität (400) fehlt, während C32 mit niedrigerer Priorität (262) enthalten ist. Dies reduziert meinen Gesamtwert erheblich."} 
	  } 
	   } 
	   } 
	 & \\ 
 

    \theutterance \stepcounter{utterance}  

    & \multicolumn{4}{p{0.6\linewidth}}{\cellcolor[rgb]{0.8,1,0.9}{%
	\makecell[{{p{\linewidth}}}]{% 
	  \tt {\tiny [A$\rangle$GM]}  
	 VORSCHLAG: {"C62", "C38", "B38", "A83", "C98", "C56", "A07", "A03", "C32"}\\ \tt ARGUMENT: {"Dieser Vorschlag enthält alle high priority Gegenstände für Sie wie A03, C56, A07 und für mich sind C62, C38, B38 von großer Wichtigkeit. Der Gesamtaufwand beträgt 2293 und ist innerhalb des maximalen Limits."} 
	  } 
	   } 
	   } 
	 & & \\ 
 

    \theutterance \stepcounter{utterance}  

    & & \multicolumn{4}{p{0.6\linewidth}}{\cellcolor[rgb]{0.9,0.9,0.9}{%
	\makecell[{{p{\linewidth}}}]{% 
	  \tt {\tiny [GM$\rangle$B]}  
	 VORSCHLAG: {"C62", "C38", "B38", "A83", "C98", "C56", "A07", "A03", "C32"}\\ \tt ARGUMENT: {"Dieser Vorschlag enthält alle high priority Gegenstände für Sie wie A03, C56, A07 und für mich sind C62, C38, B38 von großer Wichtigkeit. Der Gesamtaufwand beträgt 2293 und ist innerhalb des maximalen Limits."} 
	  } 
	   } 
	   } 
	 & \\ 
 

    \theutterance \stepcounter{utterance}  

    & & & \multicolumn{4}{p{0.6\linewidth}}{\cellcolor[rgb]{1,0.85,0.72}{%
	\makecell[{{p{\linewidth}}}]{% 
	  \tt {\tiny [GM$\langle$B]}  
	 VORSCHLAG: {"C03", "C56", "A07", "C81", "C76", "B38", "A83"}\\ \tt ARGUMENT: {"Dieser Vorschlag enthält C03 und C56 für mich mit höchster Priorität sowie B38, A83 und A07, die Sie in zahlreichen Vorschlägen als wichtig angeführt haben. Der Gesamtaufwand von 2245 ist unter der Grenze und bietet eine ausgewogene Lösung für beide Parteien."} 
	  } 
	   } 
	   } 
	  \\ 
 

    \theutterance \stepcounter{utterance}  

    & & \multicolumn{4}{p{0.6\linewidth}}{\cellcolor[rgb]{0.9,0.9,0.9}{%
	\makecell[{{p{\linewidth}}}]{% 
	  \tt {\tiny [A$\langle$GM]}  
	 VORSCHLAG: {"C03", "C56", "A07", "C81", "C76", "B38", "A83"}\\ \tt ARGUMENT: {"Dieser Vorschlag enthält C03 und C56 für mich mit höchster Priorität sowie B38, A83 und A07, die Sie in zahlreichen Vorschlägen als wichtig angeführt haben. Der Gesamtaufwand von 2245 ist unter der Grenze und bietet eine ausgewogene Lösung für beide Parteien."} 
	  } 
	   } 
	   } 
	 & \\ 
 

    \theutterance \stepcounter{utterance}  

    & \multicolumn{4}{p{0.6\linewidth}}{\cellcolor[rgb]{0.8,1,0.9}{%
	\makecell[{{p{\linewidth}}}]{% 
	  \tt {\tiny [A$\rangle$GM]}  
	 ABLEHNUNG: {"C03", "C56", "A07", "C81", "C76", "B38", "A83"}\\ \tt ARGUMENT: {"C62, C38 und C98 fehlen, die ich prioritisiere. Gegenstände wie C03 (30) und C56 (65) haben für mich nur sehr geringe Priorität, was den Gesamtnutzen stark mindert."} 
	  } 
	   } 
	   } 
	 & & \\ 
 

    \theutterance \stepcounter{utterance}  

    & & \multicolumn{4}{p{0.6\linewidth}}{\cellcolor[rgb]{0.9,0.9,0.9}{%
	\makecell[{{p{\linewidth}}}]{% 
	  \tt {\tiny [GM$\rangle$B]}  
	 ABLEHNUNG: {"C03", "C56", "A07", "C81", "C76", "B38", "A83"}\\ \tt ARGUMENT: {"C62, C38 und C98 fehlen, die ich prioritisiere. Gegenstände wie C03 (30) und C56 (65) haben für mich nur sehr geringe Priorität, was den Gesamtnutzen stark mindert."} 
	  } 
	   } 
	   } 
	 & \\ 
 

    \theutterance \stepcounter{utterance}  

    & & & \multicolumn{4}{p{0.6\linewidth}}{\cellcolor[rgb]{1,0.85,0.72}{%
	\makecell[{{p{\linewidth}}}]{% 
	  \tt {\tiny [GM$\langle$B]}  
	 VORSCHLAG: {"C03", "C56", "A03", "C62", "C38", "C98", "C76", "C32"}\\ \tt ARGUMENT: {"Der Vorschlag enthält für Sie C62, C38, C98 und A03, meine wichtigsten Gegenstände sind C03, C56 und C76. Der Gesamtaufwand von 2271 liegt unter dem Maximum."} 
	  } 
	   } 
	   } 
	  \\ 
 

    \theutterance \stepcounter{utterance}  

    & & \multicolumn{4}{p{0.6\linewidth}}{\cellcolor[rgb]{0.9,0.9,0.9}{%
	\makecell[{{p{\linewidth}}}]{% 
	  \tt {\tiny [A$\langle$GM]}  
	 VORSCHLAG: {"C03", "C56", "A03", "C62", "C38", "C98", "C76", "C32"}\\ \tt ARGUMENT: {"Der Vorschlag enthält für Sie C62, C38, C98 und A03, meine wichtigsten Gegenstände sind C03, C56 und C76. Der Gesamtaufwand von 2271 liegt unter dem Maximum."} 
	  } 
	   } 
	   } 
	 & \\ 
 

    \theutterance \stepcounter{utterance}  

    & \multicolumn{4}{p{0.6\linewidth}}{\cellcolor[rgb]{0.8,1,0.9}{%
	\makecell[{{p{\linewidth}}}]{% 
	  \tt {\tiny [A$\rangle$GM]}  
	 VORSCHLAG: {"C32", "A83", "C62", "B38", "C98", "C38", "A07", "C81", "A03", "C00"}\\ \tt ARGUMENT: {"Dieser Vorschlag enthält Gegenstände wie C32, A83, C62, B38 und C98, die ich prioritisiere, sowie A07, C81, C00, A03, die auch in Ihren bisherigen Vorschlägen vorkommen. Der Gesamtaufwand liegt bei 2268 und ist unterhalb des Maximums."} 
	  } 
	   } 
	   } 
	 & & \\ 
 

    \theutterance \stepcounter{utterance}  

    & & \multicolumn{4}{p{0.6\linewidth}}{\cellcolor[rgb]{0.9,0.9,0.9}{%
	\makecell[{{p{\linewidth}}}]{% 
	  \tt {\tiny [GM$\rangle$B]}  
	 VORSCHLAG: {"C32", "A83", "C62", "B38", "C98", "C38", "A07", "C81", "A03", "C00"}\\ \tt ARGUMENT: {"Dieser Vorschlag enthält Gegenstände wie C32, A83, C62, B38 und C98, die ich prioritisiere, sowie A07, C81, C00, A03, die auch in Ihren bisherigen Vorschlägen vorkommen. Der Gesamtaufwand liegt bei 2268 und ist unterhalb des Maximums."} 
	  } 
	   } 
	   } 
	 & \\ 
 

    \theutterance \stepcounter{utterance}  

    & & & \multicolumn{4}{p{0.6\linewidth}}{\cellcolor[rgb]{1,0.85,0.72}{%
	\makecell[{{p{\linewidth}}}]{% 
	  \tt {\tiny [GM$\langle$B]}  
	 VORSCHLAG: {"C03", "C56", "C81", "A07", "C76", "C10", "A03"}\\ \tt ARGUMENT: {"Dieser Vorschlag enthält meine höchsten Prioritäten C03 und C56 sowie A07, C81, C76 und C10, die für Sie bereits priorisiert wurden. Der Gesamtaufwand von 2257 liegt unter dem Limit und sichert einen fairen Nutzen für beide Parteien."} 
	  } 
	   } 
	   } 
	  \\ 
 

    \theutterance \stepcounter{utterance}  

    & & \multicolumn{4}{p{0.6\linewidth}}{\cellcolor[rgb]{0.9,0.9,0.9}{%
	\makecell[{{p{\linewidth}}}]{% 
	  \tt {\tiny [A$\langle$GM]}  
	 VORSCHLAG: {"C03", "C56", "C81", "A07", "C76", "C10", "A03"}\\ \tt ARGUMENT: {"Dieser Vorschlag enthält meine höchsten Prioritäten C03 und C56 sowie A07, C81, C76 und C10, die für Sie bereits priorisiert wurden. Der Gesamtaufwand von 2257 liegt unter dem Limit und sichert einen fairen Nutzen für beide Parteien."} 
	  } 
	   } 
	   } 
	 & \\ 
 

    \theutterance \stepcounter{utterance}  

    & \multicolumn{4}{p{0.6\linewidth}}{\cellcolor[rgb]{0.8,1,0.9}{%
	\makecell[{{p{\linewidth}}}]{% 
	  \tt {\tiny [A$\rangle$GM]}  
	 ABLEHNUNG: {"C03", "C56", "C81", "A07", "C76", "C10", "A03"}\\ \tt ARGUMENT: {"Die Gegenstände C03 und C56 haben für mich sehr niedrige Priorität, während wichtige Items wie C62, C38, B38, C98 und A83 fehlen, die deutlich höheren Nutzen hätten."} 
	  } 
	   } 
	   } 
	 & & \\ 
 

    \theutterance \stepcounter{utterance}  

    & & \multicolumn{4}{p{0.6\linewidth}}{\cellcolor[rgb]{0.9,0.9,0.9}{%
	\makecell[{{p{\linewidth}}}]{% 
	  \tt {\tiny [GM$\rangle$B]}  
	 ABLEHNUNG: {"C03", "C56", "C81", "A07", "C76", "C10", "A03"}\\ \tt ARGUMENT: {"Die Gegenstände C03 und C56 haben für mich sehr niedrige Priorität, während wichtige Items wie C62, C38, B38, C98 und A83 fehlen, die deutlich höheren Nutzen hätten."} 
	  } 
	   } 
	   } 
	 & \\ 
 

    \theutterance \stepcounter{utterance}  

    & & & \multicolumn{4}{p{0.6\linewidth}}{\cellcolor[rgb]{1,0.85,0.72}{%
	\makecell[{{p{\linewidth}}}]{% 
	  \tt {\tiny [GM$\langle$B]}  
	 VORSCHLAG: {"C62", "C38", "B38", "C98", "A83", "C76", "A07", "C81", "A03", "C32"}\\ \tt ARGUMENT: {"Dieser Vorschlag beinhaltet alle von Ihnen gewünschten Hochprioritäts-Gegenstände wie C62, C38, B38, C98 und A83 sowie für mich wesentliche Items wie C76, C81, A07, A03 und C32. Der Gesamtaufwand von 2107 liegt deutlich unter der Grenze und sichert einen optimalen Nutzen für beide Seiten."} 
	  } 
	   } 
	   } 
	  \\ 
 

    \theutterance \stepcounter{utterance}  

    & & \multicolumn{4}{p{0.6\linewidth}}{\cellcolor[rgb]{0.9,0.9,0.9}{%
	\makecell[{{p{\linewidth}}}]{% 
	  \tt {\tiny [A$\langle$GM]}  
	 VORSCHLAG: {"C62", "C38", "B38", "C98", "A83", "C76", "A07", "C81", "A03", "C32"}\\ \tt ARGUMENT: {"Dieser Vorschlag beinhaltet alle von Ihnen gewünschten Hochprioritäts-Gegenstände wie C62, C38, B38, C98 und A83 sowie für mich wesentliche Items wie C76, C81, A07, A03 und C32. Der Gesamtaufwand von 2107 liegt deutlich unter der Grenze und sichert einen optimalen Nutzen für beide Seiten."} 
	  } 
	   } 
	   } 
	 & \\ 
 

    \theutterance \stepcounter{utterance}  

    & \multicolumn{4}{p{0.6\linewidth}}{\cellcolor[rgb]{0.8,1,0.9}{%
	\makecell[{{p{\linewidth}}}]{% 
	  \tt {\tiny [A$\rangle$GM]}  
	 ZUSTIMMUNG: {"C62", "C38", "B38", "C98", "A83", "C76", "A07", "C81", "A03", "C32"} 
	  } 
	   } 
	   } 
	 & & \\ 
 

    \theutterance \stepcounter{utterance}  

    & & \multicolumn{4}{p{0.6\linewidth}}{\cellcolor[rgb]{0.9,0.9,0.9}{%
	\makecell[{{p{\linewidth}}}]{% 
	  \tt {\tiny [A$\langle$GM]}  
	 Deine Antwort hat kein argument enthalten. Du muss das strukturierte Format ARGUMENT: {...} irgendwo in theiner Antwort verwenden. Versuche es erneut 
	  } 
	   } 
	   } 
	 & \\ 
 

    \theutterance \stepcounter{utterance}  

    & \multicolumn{4}{p{0.6\linewidth}}{\cellcolor[rgb]{0.8,1,0.9}{%
	\makecell[{{p{\linewidth}}}]{% 
	  \tt {\tiny [A$\rangle$GM]}  
	 ZUSTIMMUNG: {"C62", "C38", "B38", "C98", "A83", "C76", "A07", "C81", "A03", "C32"}\\ \tt ARGUMENT: {"Der Vorschlag berücksichtigt alle high priority Gegenstände für beide Parteien und liegt innerhalb des Maximalaufwands. Eine faire Einigung ist nun erreicht."} 
	  } 
	   } 
	   } 
	 & & \\ 
 

    \theutterance \stepcounter{utterance}  

    & & & \multicolumn{2}{p{0.3\linewidth}}{\cellcolor[rgb]{0.95,0.95,0.95}{%
	\makecell[{{p{\linewidth}}}]{% 
	  \tt {\tiny [GM$|$GM]}  
	 game successful 
	  } 
	   } 
	   } 
	 & & \\ 
 

    \theutterance \stepcounter{utterance}  

    & & & \multicolumn{2}{p{0.3\linewidth}}{\cellcolor[rgb]{0.95,0.95,0.95}{%
	\makecell[{{p{\linewidth}}}]{% 
	  \tt {\tiny [GM$|$GM]}  
	 end game 
	  } 
	   } 
	   } 
	 & & \\ 
 

\end{supertabular}
}

\end{document}
